\chapter{Code and reproducibility}\label{app:code}

The scripts, the data files they write and the reruns made for the book are not printed here: they are released with the PDF, in the book's repository, \url{https://github.com/VadimSokolov/computing_rh}, beside the \LaTeX\ source of the book. Most scripts are in the folder \texttt{code/} and the data files they write in \texttt{data/}, the reruns and reconstructions made for the book in \texttt{repro/}, and the scripts of some chapters in other folders. The file \texttt{code/\allowbreak README.md} lists every script with its purpose, the data files it reads and writes, the results of the book that it produces and its usage, and describes the reruns of September 2026, the computations that have no script in the archive and the gaps that remain.

The scripts run with Python 3.11 and the libraries of \texttt{requirements.txt}: \texttt{python-flint} 0.9.0, the Python interface to FLINT and Arb \cite{Arb}, for ball arithmetic and the certificates, and \texttt{mpmath} \cite{mpmath} for reference values. The long computations ran on the Hopper cluster of George Mason University. The scripts of the later chapters keep a log of each run beside the data file, those of Chapters~\ref{chap:yoshida} to \ref{chap:gl} take a \texttt{--quick} option that runs a reduced case in seconds, and where a chapter printed a value that was recomputed for the book, the data file or the log records the value computed beside it and says whether the two agree. Where the book says that a statement is proved in ball arithmetic, the script named beside it is the certificate, and its output records the enclosures.

\section{The P\'olya evaluators}\label{app:code:eval}

The evaluator of Section~\ref{ch:polya} on the real $u$ axis is \texttt{code/polya\_flint.py}. In ball arithmetic it computes the node values $\Phi_j=\Phi(jh)$, $0\le j\le J$, from the theta series \eqref{eq:Phi}, and from them $P(E)+P(E^{-1})$ and $Q(E)-Q(E^{-1})$ with $E=e^{sh}$ and the polynomials $P$ and $Q$ of Section~\ref{sec:ch4:trapezoid}, the trapezoid approximations to $\xi(\frac12+s)$ and $\xi'(\frac12+s)$. The defaults are $240$ bits, $h=1/200$ and $|u|\le3$, so that $J=600$. Most runs use $420$ bits, the parameters of Section~\ref{sec:ch4:trapezoid}; the Keiper Li coefficients of Section~\ref{ch:li} use $600$ bits, and the Selberg computation of Section~\ref{ch:sin2more} uses $300$.

The tilted evaluator \texttt{code/tilted.py} implements Proposition~\ref{prop:tilt}. With the nodes $u_j=x_j+i\eta$, $x_j=jh$, $|j|\le J$, where $J=\lceil x_{\max}/h\rceil$ for the cut $|x|\le x_{\max}$, it returns the trapezoid sums of the two integrals of the proposition. The production scripts evaluate each exponential directly, not by Horner's rule, whose ball radii grow as described in Section~\ref{sec:ch4:balls}. The Lehmer pair and its de Bruijn Newman flow use $\delta=0.01$, $256$ bits, $h=6\times10^{-4}$ and $|x|\le3.9$; the count at height $1000$ in Section~\ref{sec:ch5:count} uses $\delta=0.1$, $200$ bits, $h=0.004$ and $|x|\le4.2$; the Laplace transform of the clock for Sections~\ref{ch:clocklaw} and \ref{ch:boundaries} switches to the tilted evaluator with $\delta=0.1$, $320$ bits, $h=0.003$ and $|x|\le4.5$ when $|\im s|\ge120$; and the runs of Table~\ref{tab:tilt} take $\delta$, the precision, $h$, the cut and the heights from the command line, which \texttt{code/\allowbreak README.md} records.

In both evaluators the radii of the returned balls bound the rounding errors of the finite sums, and nothing else. Three errors lie outside them: the discretisation error of the trapezoid rule, bounded in Section~\ref{sec:ch4:trapezoid}; the cut of the axis at $|u|\le3$, where $\Phi<10^{-500}$, or at $|x|\le x_{\max}$ on the tilted line; and the truncation of the theta series, whose loop stops at the first term with midpoint below $2^{-b-20}$ at working precision $b$ and does not add the tail to the radius. The stopping test compares double precision numbers, so for $b\ge1055$ its threshold underflows to zero and the test is never met. The untilted route needs more than $1054$ bits for thirty digits above height about $840$ (Table~\ref{tab:ch4:budget}), and there the test has to be made in ball arithmetic. The runs in this book use at most $600$ bits and are not affected.
